\section{The Analytical Framework: Formalizing the
	Transformation Elasticity}

\paragraph{3.0}
This section translates the distributive--technological
mechanism of Section~2 into a formal model that yields
estimable restrictions. \citet{CajasGuajirro2024} first
advanced the cost-minimization problem over mechanization
growth; we extend it through three subsections of
increasing complexity. Section~3.1 formalizes aggregate
capital-deepening under Leontief technology and derives
the transformation elasticity $\theta$ endogenously from
the cost-minimizing choice of technique. Section~3.2
unbundles aggregate capital into extensive nonresidential
structures and intensive machinery, and derives the
state-dependent mechanization gap $\kappa_t$ under a
dual-channel viability condition. Section~3.3 imposes the
peripheral foreign-exchange bottleneck as a binding
constraint on $\kappa_t$, generating a corner solution
that truncates $\theta$. Each subsection nests the
preceding one: setting the heterogeneity parameter to
zero recovers the baseline model; relaxing the
foreign-exchange constraint recovers the center model.
The empirical analysis bounds this multi-sector mechanism
to a single-sector aggregate, which masks internal
sectoral divergences but preserves the core transformation
logic. Section~3.4 synthesizes the theoretical claims of
the three models. Section~4 consolidates the
identification strategy and derives the cointegrating
specifications.

\subsection{The Baseline Model of Cost-Minimizing
	Mechanization (Homogeneous Capital Stock)}

\paragraph{3.1}
Actual aggregate output is represented by a Leontief
production function with fixed technical coefficients:
\begin{equation}
	Y = \min\left( A L,\; \mu B K \right)
	\label{eq:31:leontief}
\end{equation}
where $A = Y/L$ is actual labor productivity, $\mu = Y/Y^p$
is capacity utilization, and $B = Y^p/K$ is maximum capital
productivity. This production structure governs the rate of
profit through the fundamental heterodox accounting identity:
\begin{equation}
	r = \mu B (1 - \omega)
	\label{eq:31:profit_identity}
\end{equation}
where $\omega$ is the wage share. Standard macroeconomic
models treat $B$ as an exogenous technological given and $\mu$
as a behavioral constant smoothly targeted by a unified firm.
We reject both assumptions. The components of
\eqref{eq:31:profit_identity} are endogenous outcomes of
class struggle and decentralized market anarchy, as
established in Section~2. This subsection formalizes how
distributive pressure forces cost-minimizing firms to alter
their choice of technique, endogenously determining $B$.

\subsubsection{Tactical Choice of Technique and Induced
	Innovation}

\paragraph{3.2}
Firms select a mechanization growth rate $q$ to maximize the
net productivity surplus. Let $a$ denote the growth rate of
labor productivity. The effective distributive cost of
deploying mechanization growth is $q(1-\omega)$. The firm
solves:
\begin{equation}
	\max_{q}\; c(q) = a - q(1-\omega)
	\label{eq:31:objective}
\end{equation}
This objective captures the induced-innovation logic: firms
accelerate mechanization to expand labor productivity growth,
but the distributive cost of capital intensity holds them
back. A technological frontier bounds this choice:
\begin{equation}
	a = \Phi(q) = \lambda_0 + \lambda_1 q + \lambda_2 q^2
	\label{eq:31:frontier}
\end{equation}

\paragraph{3.3}
\textbf{Assumption 1.} \textit{Concavity of the
	mechanization frontier.} The frontier parameters satisfy
$\lambda_0 > 0$, $\lambda_1 > 0$, and $\lambda_2 < 0$. The
base productivity growth $\lambda_0$ sets the floor without
advanced machinery. The linear coefficient $\lambda_1$
captures the direct productivity payoff of mechanization. The
strict concavity $\lambda_2 < 0$ ensures diminishing marginal
returns to mechanization.

\paragraph{3.4}
Substituting \eqref{eq:31:frontier} into
\eqref{eq:31:objective} and solving the first-order condition
yields the optimal mechanization growth:
\begin{align}
	c(q) &= \lambda_0 + \left[\lambda_1 - (1-\omega)\right] q
	+ \lambda_2 q^2 \label{eq:31:unconstrained} \\
	\frac{dc}{dq} &= \lambda_1 - (1-\omega) + 2\lambda_2 q^*
	= 0 \label{eq:31:foc} \\
	q^*(\omega) &= \frac{(1-\omega) - \lambda_1}{2\lambda_2}
	\label{eq:31:qstar}
\end{align}
The strict concavity ($c''(q) = 2\lambda_2 < 0$) guarantees
a unique maximum.

\paragraph{3.5}
\textbf{Proposition 1.} \textit{Induced innovation.} Under
Assumption~1, the optimal mechanization growth is strictly
increasing in the wage share:
\begin{equation}
	\frac{\partial q^*}{\partial \omega}
	= -\frac{1}{2\lambda_2} > 0
	\label{eq:31:comparative_static}
\end{equation}
A rise in the wage share forces cost-minimizing firms to
substitute machinery for labor. Rising wages drive
mechanization directly through the first-order condition,
pushing optimal capital intensity growth upward.
\textit{Proof.} Immediate from differentiating
\eqref{eq:31:qstar}. $\hfill\blacksquare$

\paragraph{3.X}
\textbf{Remark 1.} The objective
\eqref{eq:31:objective} does not assert that the firm is a
coherent optimizer. Following \citet{Vidal2022}, the firm is
internally divided between the \textit{control imperative}
(valorization: cost reduction through mechanization) and the
\textit{cooperation imperative} (production: the need for
worker effort, tacit knowledge, and problem-solving
capacity). The formal model captures the cost-minimizing
pole---the Okishio viability logic---while the constraints on
the optimization problem capture the limits imposed by the
cooperation imperative and by external conditions. The firm
maximizes net productivity surplus---labor productivity
growth minus the distributive cost of mechanization---subject
to diminishing returns. The firm's actual mechanization
choice is a contested settlement between these imperatives,
not a coherent optimum. What the model demonstrates is the
cost-reducing direction in which distributive pressure pushes
the settlement, not the settlement itself.

\subsubsection{Endogenous Transformation Elasticity}

\paragraph{3.6}
Substituting the optimal technique back into the
technological frontier recovers the normal growth rate of
labor productivity under smooth reproduction:
\begin{equation}
	a^N(\omega) = \lambda_0 + \lambda_1 q^*(\omega)
	+ \lambda_2 \left[ q^*(\omega) \right]^2
	\label{eq:31:aN}
\end{equation}
This optimal technique dictates the structural trajectory of
the Leontief parameter $B$. By definition, maximum capital
productivity is $B = Y^p / K$. In dynamic growth rates, the
growth of capital productivity is:
\begin{equation}
	b \equiv g_{Y^p} - g_K = a - q
	\label{eq:31:b_def}
\end{equation}
Evaluated at the smooth reproduction benchmark:
$b^N(\omega) = a^N(\omega) - q^*(\omega)$.

\paragraph{3.7}
\textbf{Definition 1.} \textit{Transformation elasticity.}
The transformation elasticity $\theta$ is the structural
parameter that translates capital accumulation into productive
capacity growth:
\begin{equation}
	g_{Y^p} = \theta\, g_K
	\label{eq:31:theta_def}
\end{equation}
Substituting \eqref{eq:31:theta_def} into
\eqref{eq:31:b_def} reveals the link between the level
coefficient $B$ and the elasticity $\theta$:
\begin{equation}
	\theta = 1 + \frac{b}{g_K}
	\label{eq:31:theta_link}
\end{equation}

\paragraph{3.8}
To evaluate this elasticity, we anchor the optimal technique
to a baseline rate of capital accumulation. Following the
normal reproduction path, the normal accumulation rate is
driven by the normal profit rate and the normal
recapitalization share, net of depreciation:
\begin{equation}
	g_K^N = \chi^N r^N - \delta
	\label{eq:31:gKN}
\end{equation}
Expressed dynamically relative to this normal path, the
aggregate transformation elasticity becomes:
\begin{equation}
	\theta(\omega) = 1 + \frac{b^N(\omega)}{g_K^N}
	= 1 + \frac{a^N(\omega) - q^*(\omega)}{g_K^N}
	\label{eq:31:theta_omega}
\end{equation}

\paragraph{3.9}
\textbf{Proposition 2.} \textit{Endogeneity of the
	transformation elasticity.} Under Assumption~1, the
transformation elasticity $\theta$ is not an exogenous
technical constant. It emerges endogenously from the
cost-minimizing choice of technique. When distributive
pressure forces firms to alter mechanization via
\eqref{eq:31:comparative_static}, it simultaneously shifts
the structural efficiency of accumulation via
\eqref{eq:31:theta_omega}. The elasticity $\theta$ is
therefore a function of the wage share $\omega$, the frontier
parameters $\lambda_0, \lambda_1, \lambda_2$, and the normal
accumulation rate $g_K^N$.


\subsection{Heterogeneous Capital, the Mechanization Gap,
	and State-Dependent Viability}

\paragraph{3.10}
Mainstream macroeconomics treats the capital stock as a
homogeneous aggregate, abstracting from the internal
architecture of the firm. Following the theory of the
divided firm \citep{Vidal2015, Vidal2019}, we unbundle
aggregate productive capital into extensive nonresidential
structures ($K^{NRC}$) and intensive machinery ($K^{ME}$).
These components do not enter the labor process
symmetrically. Structures form the passive infrastructure
envelope---the factory floor and logistics shell. Machinery
is the active component that restructures the production of
relative surplus value \citep[Vol.~I, Ch.~15]{MarxKVI}.

\paragraph{3.11}
\textbf{Assumption 2.} \textit{Heterogeneous capital
	decomposition.} Aggregate productive capital decomposes as
$K^{cap}_t = K^{NRC}_t + K^{ME}_t$. Nonresidential
structures evolve independently of short-run distributive
conflict ($\omega$) and follow a normal accumulation
corridor governed by long-run reproduction conditions:
\begin{equation}
	g_{K^{NRC}, t}^N = \chi^N r^N_t - \delta
	\label{eq:32:gKNRC}
\end{equation}
This extensive envelope $\Gamma(K^{NRC})$ bounds the space
of tactical technique choice.

\subsubsection{The Extensive Envelope and the Intensive
	Margin}

\paragraph{3.12}
Given the fixed infrastructural envelope, the
\textbf{intensive margin} of accumulation is the tactical
choice of the mechanization gap, defined in natural logs as:
\begin{equation}
	\kappa_t = k^{ME}_t - k^{NRC}_t
	\label{eq:32:kappa_def}
\end{equation}
This gap measures mechanization intensity relative to the
extensive plant scale. The capacity identity decomposes
potential output into the extensive envelope and the
intensive capacity yield $\Psi_t$:
\begin{equation}
	y^p_t = \theta_{NRC} k^{NRC}_t
	+ \Psi_t(\kappa_t, \omega_{t-1}, p^\kappa_t)
	\label{eq:32:capacity_decomp}
\end{equation}
The intensive yield $\Psi_t$ is not a fixed engineering
constant. It is the outcome of a decentralized
cost-minimization problem governed by the regime of
accumulation, as formalized below.

\subsubsection{Cost-Surplus Maximization and the Okishio
	Viability Yield}

\paragraph{3.13}
The firm chooses the mechanization gap $\kappa_t$ to
maximize the net viability surplus, subject to the
regime-conditioned mechanization frontier:
\begin{equation}
	\Phi(\kappa_t) = \lambda_0 + \lambda_1 \kappa_t
	+ \lambda_2 \kappa_t^2
	\label{eq:32:frontier_kappa}
\end{equation}
The \textit{realized} capacity yield is governed by the
\textbf{Marx--Okishio viability condition}: a
cost-minimizing firm deploys a new technique only if it
reduces unit costs at current distributive and price
parameters \citep{Okishio1961, Basu2022}. We formalize
this as a cost-surplus maximization problem:
\begin{equation}
	\max_{\kappa_t}\; c(\kappa_t)
	= \Phi(\kappa_t) - \kappa_t \cdot MC(\omega_{t-1}, p^\kappa_t)
	\label{eq:32:objective_kappa}
\end{equation}

\paragraph{3.14}
The marginal viability cost $MC(\omega_{t-1}, p^\kappa_t)$
is determined by two channels:
\begin{enumerate}
	\item \textbf{The distributive dividend (labor saving).}
	Increasing $\kappa_t$ substitutes machinery for labor. The
	value of this displacement is proportional to the lagged
	wage share ($\omega_{t-1}$), reflecting the wage
	settlement period prior to current capacity realization.
	A higher wage share \textit{reduces} the net cost of
	mechanization.
	\item \textbf{The relative price penalty (capital cost).}
	Increasing $\kappa_t$ requires accumulating machinery
	faster than structures. The cost of this composition shift
	is proportional to the log relative price of machinery
	($p^\kappa_t = \ln P^{ME}_t - \ln P^{NRC}_t$). A higher
	relative price \textit{increases} the marginal cost.
\end{enumerate}
We specify the marginal viability cost as:
\begin{equation}
	MC(\omega_{t-1}, p^\kappa_t)
	= \gamma_0 - \gamma_1 \omega_{t-1} + \gamma_2 p^\kappa_t
	\label{eq:32:MC}
\end{equation}
where $\gamma_1 > 0$ captures the wage-savings dividend and
$\gamma_2 > 0$ captures the relative price penalty.

\paragraph{3.15}
Substituting \eqref{eq:32:frontier_kappa} and
\eqref{eq:32:MC} into \eqref{eq:32:objective_kappa} and
solving the first-order condition yields the optimal
mechanization gap:
\begin{align}
	c(\kappa_t) &= \lambda_0
	+ \left[\lambda_1 - (\gamma_0 - \gamma_1 \omega_{t-1}
	+ \gamma_2 p^\kappa_t)\right] \kappa_t
	+ \lambda_2 \kappa_t^2
	\label{eq:32:unconstrained_kappa} \\
	\frac{dc}{d\kappa_t} &= \lambda_1 - \gamma_0
	+ \gamma_1 \omega_{t-1} - \gamma_2 p^\kappa_t
	+ 2\lambda_2 \kappa_t^* = 0
	\label{eq:32:foc_kappa} \\
	\kappa_t^*(\omega_{t-1}, p^\kappa_t)
	&= \frac{(\gamma_0 - \lambda_1) - \gamma_1 \omega_{t-1}
		+ \gamma_2 p^\kappa_t}{2\lambda_2}
	\label{eq:32:kappa_star}
\end{align}

\paragraph{3.16}
\textbf{Proposition 3.} \textit{Dual-channel induced
	innovation.} Under Assumption~1 (concavity) and
Assumption~2 (heterogeneous capital), the optimal
mechanization gap satisfies:
\begin{align}
	\frac{\partial \kappa_t^*}{\partial \omega_{t-1}}
	&= -\frac{\gamma_1}{2\lambda_2} > 0
	\label{eq:32:cs_omega} \\
	\frac{\partial \kappa_t^*}{\partial p^\kappa_t}
	&= \frac{\gamma_2}{2\lambda_2} < 0
	\label{eq:32:cs_price}
\end{align}
A rising wage share forces cost-minimizing firms to widen
the mechanization gap, substituting machinery for labor.
A rising relative price of machinery suppresses the optimal
gap, as the structural cost of the composition shift
penalizes mechanization. \textit{Proof.} Immediate from
differentiating \eqref{eq:32:kappa_star}. $\hfill\blacksquare$

\subsubsection{The Marginal Capacity Payoff and the
	Econometric Bridge}

\paragraph{3.17}
The macroeconomic translation from the mechanization gap to
potential capacity is given by the \textbf{marginal
	capacity payoff} ($MP_\kappa$), the derivative of the
intensive capacity yield $\Psi_t$ with respect to $\kappa_t$.
By the envelope theorem:
\begin{align}
	MP_\kappa &= \frac{\partial \Psi_t}{\partial \kappa_t}
	= \frac{\partial \Phi}{\partial \kappa_t}
	- MC(\omega_{t-1}, p^\kappa_t)
	\label{eq:32:MP_env} \\
	&= (\lambda_1 + 2\lambda_2 \kappa_t)
	- (\gamma_0 - \gamma_1 \omega_{t-1} + \gamma_2 p^\kappa_t)
	\label{eq:32:MP_expanded} \\
	&= \underbrace{(\lambda_1 - \gamma_0
		+ 2\lambda_2 \kappa_t)}_{\text{autonomous yield}}
	+ \gamma_1 \omega_{t-1} - \gamma_2 p^\kappa_t
	\label{eq:32:MP_rearranged}
\end{align}

\paragraph{3.18}
Integrating the marginal payoff with respect to $\kappa_t$
(treating $\omega_{t-1}$ and $p^\kappa_t$ as predetermined
shifters) and defining $\theta_\kappa$,
$\theta_{\omega\kappa} \equiv \gamma_1$, and
$\theta_{\kappa p\kappa} \equiv -\gamma_2$ yields the
intensive capacity yield function:
\begin{equation}
	\Psi_t = \theta_\kappa \kappa_t
	+ \theta_{\omega\kappa} (\omega_{t-1} \kappa_t)
	+ \theta_{\kappa p\kappa} (p^\kappa_t \kappa_t)
	\label{eq:32:Psi}
\end{equation}
Substituting \eqref{eq:32:Psi} back into
\eqref{eq:32:capacity_decomp} recovers the structural
capacity equation:
\begin{equation}
	y^p_t = \theta_{NRC} k^{NRC}_t + \theta_\kappa \kappa_t
	+ \theta_{\omega\kappa} (\omega_{t-1} \kappa_t)
	+ \theta_{\kappa p\kappa} (p^\kappa_t \kappa_t)
	\label{eq:32:capacity_eq}
\end{equation}
The interaction terms are not ad-hoc econometric additions.
They are the macroeconomic footprint of the divided firm's
Okishio viability constraint. Distributive pressure and
relative price enter because they alter the marginal
capacity payoff of the mechanization gap.

\subsubsection{The Aggregate Transformation Elasticity}

\paragraph{3.19}
The aggregate transformation elasticity
($\theta^{SpecB}_t$) is the ratio of the total differentials
of potential capacity and aggregate log capital along the
historical accumulation path. The total differential of
\eqref{eq:32:capacity_eq}, holding the slow-moving state
variables fixed, separates the extensive envelope from the
intensive margin:
\begin{align}
	dy^p_t &= \theta_{NRC} dk^{NRC}_t + MP_\kappa d\kappa_t
	\label{eq:32:dy_total} \\
	&= \theta_{NRC} dk^{NRC}_t
	+ \left(\theta_\kappa + \theta_{\omega\kappa} \omega_{t-1}
	+ \theta_{\kappa p\kappa} p^\kappa_t\right) d\kappa_t
	\label{eq:32:dy_expanded}
\end{align}

\paragraph{3.20}
The variation in aggregate capital follows from the Divisia
index of the heterogeneous components. Since
$s_{ME,t} + s_{NRC,t} = 1$ and
$d\kappa_t = dk^{ME}_t - dk^{NRC}_t$:
\begin{align}
	dk^{cap}_t &\approx s_{ME,t} dk^{ME}_t
	+ s_{NRC,t} dk^{NRC}_t
	\label{eq:32:dk_divisia} \\
	&= dk^{NRC}_t + s_{ME,t} d\kappa_t
	\label{eq:32:dk_simplified}
\end{align}
Dividing \eqref{eq:32:dy_expanded} by
\eqref{eq:32:dk_simplified} yields the elasticity along the
empirical growth path:
\begin{equation}
	\theta^{SpecB}_t
	= \frac{\theta_{NRC}
		+ \left(\theta_\kappa + \theta_{\omega\kappa} \omega_{t-1}
		+ \theta_{\kappa p\kappa} p^\kappa_t\right)
		\frac{d\kappa_t}{dk^{NRC}_t}}
	{1 + s_{ME,t} \frac{d\kappa_t}{dk^{NRC}_t}}
	\label{eq:32:theta_SpecB}
\end{equation}

\paragraph{3.21}
The aggregate transformation elasticity
$\theta^{SpecB}_t$ in \eqref{eq:32:theta_SpecB} is
state-dependent: its value shifts with the distributive
channel ($\omega_{t-1}$), the price channel
($p^\kappa_t$), and the empirical growth path
$d\kappa_t / dk^{NRC}_t$. Whether $\theta^{SpecB}_t$
exceeds or falls below unity is an empirical question
determined by the interaction of these channels under a
given accumulation regime. The identification strategy in
Section~3.4 specifies how this regime-dependence is
estimated and tested.



\subsection{The Peripheral Extension: Balance-of-Payments
	Constraint and the Corner Solution}

\paragraph{3.22}
In the capitalist periphery, the intensive margin of
accumulation is structurally tethered to the world market and
constrained by domestic social friction. Following
\citet{Prebisch1981} and \citet{Kaldor1959}, peripheral
capitalism is characterized by high consumption propensities
out of profits, structural inelasticities in agriculture, and
a chronic vulnerability to balance-of-payments crises
\citep{Ahumada2019}. This introduces a dual realization
bottleneck: the macroeconomic capacity to import and the
inelastic domestic supply of wage-goods. Section~2.5
established this constraint qualitatively as the peripheral
accumulation wedge. This subsection formalizes it as a strict
macroeconomic boundary constraint on the cost-minimizing
firm's choice of technique, integrating the structuralist
insights of \citet{Prebisch1981} into the heterogeneous
capital framework of Section~3.2.

\subsubsection{The Import-Biased Intensive Margin and the
	Balance of Payments Constraint}

\paragraph{3.23}
In Section~3.2, the cost-minimizing firm chooses the
mechanization gap $\kappa_t$ to maximize the net viability
surplus:
\begin{equation}
	\max_{\kappa_t}\; c(\kappa_t)
	= \Phi(\kappa_t) - \kappa_t \cdot MC(\omega_{t-1}, p^\kappa_t)
	\label{eq:33:objective_periphery}
\end{equation}
where $\Phi(\kappa_t) = \lambda_0 + \lambda_1 \kappa_t
+ \lambda_2 \kappa_t^2$ is the regime-conditioned
mechanization frontier, and $MC(\omega_{t-1}, p^\kappa_t)
= \gamma_0 - \gamma_1 \omega_{t-1} + \gamma_2 p^\kappa_t$
is the marginal viability cost of the gap. In the core, this
optimization proceeds unconstrained, yielding the optimal
mechanization gap $\kappa_t^*$. In the periphery, however,
active mechanization is strictly import-biased: widening the
mechanization gap requires imported capital goods, specialized
machinery, and technologically non-substitutable spare parts \citep{Vernengo2006}.

\paragraph{3.24}
\textbf{Assumption 3.} \textit{Import-biased mechanization.}
The physical import requirement per unit of output
($m^{req}_t$) is a linear function of the chosen technique:
\begin{equation}
	m^{req}_t = \nu \cdot \kappa_t
	\label{eq:33:import_req}
\end{equation}
where $\nu > 0$ captures the structural import-intensity of
the mechanization gap. Simultaneously, the physical capacity
to satisfy this requirement is bounded by the macroeconomic
\textbf{capacity to import} per unit of output
($m^{avail}_t = M^{avail}_t / Y_t$). Following
\citet{Prebisch1981}, peripheral export growth is governed by
the center's income elasticity of demand for primary goods
($\epsilon_{m,c} < 1$), structurally lagging behind domestic
import demand. This external asymmetry, compounded by high
consumption propensities out of capitalist profits
\citep{Kaldor1959}, severely squeezes the foreign exchange
available for productive accumulation. This imposes a strict
external ceiling on the choice of technique:
\begin{equation}
	\nu \cdot \kappa_t \le m^{avail}_t
	\implies \kappa_t \le \bar{\kappa}_t
	\equiv \frac{m^{avail}_t}{\nu}
	\label{eq:33:bop_ceiling}
\end{equation}

\subsubsection{The Macroeconomic Allocation of Foreign
	Exchange}

\paragraph{3.25}
\textbf{Assumption 4.} \textit{Balance-of-payments
	constraint.} Total import demand cannot exceed export
capacity. The total supply of foreign exchange per unit of
output is driven by primary exports
($x_t = p_t^x X_t / Y_t$). Due to the center's low income
elasticity for primary goods, $x_t$ is structurally rigid.
The total demand for foreign exchange per unit of output
consists of three competing components:
\begin{enumerate}
	\item \textbf{Wage-goods imports} ($m^W_t$): driven by the
	wage share $\omega_{t-1}$ and the structural inelasticity
	of domestic agriculture \citep{Kaldor1959}.
	$m^W_t = m_w \omega_{t-1}$.
	\item \textbf{Conspicuous consumption imports} ($m^C_t$):
	driven by the profit share $\pi_t$ and the high propensity
	of the peripheral upper strata to imitate the consumption
	patterns of the center \citep{Prebisch1981, Kaldor1959}.
	$m^C_t = m_c c_k \pi_t$.
	\item \textbf{Productive capital imports} ($m^K_t$): driven
	by the mechanization gap. $m^K_t = \nu \kappa_t$.
\end{enumerate}

\paragraph{3.26}
\textbf{Assumption 5.} \textit{Surplus leakage.} The
propensity to consume out of profits satisfies
$c_k > c_k^{\min}$, where $c_k^{\min}$ is the threshold
below which domestic savings suffice to finance capital
formation without external borrowing. Under the
developmentalist regime, the leakage channel is conspicuous
consumption imports by domestic elites
\citep{Prebisch1981, Kaldor1959}. Under the neoliberal
regime, the leakage channel is profit repatriation by
multinational corporations and capital flight
\citep{Ahumada2019}. In both cases, $c_k > c_k^{\min}$
ensures that the effective profit share available for
productive reinvestment is compressed below the nominal
level: $\pi_t^{eff} = \pi_t \cdot \phi_t < \pi_t$.

\paragraph{3.27}
The balance of payments constraint requires:
\begin{equation}
	m_w \omega_{t-1} + m_c c_k \pi_t + \nu \kappa_t \le x_t
	\label{eq:33:bop_full}
\end{equation}
Using the distributive identity $\pi_t = 1 - \omega_{t-1}$,
we substitute and solve for the maximum feasible mechanization
gap:
\begin{align}
	m_w \omega_{t-1} + m_c c_k (1 - \omega_{t-1})
	+ \nu \kappa_t &\le x_t
	\label{eq:33:bop_sub} \\
	\nu \kappa_t &\le x_t - m_c c_k
	+ (m_c c_k - m_w) \omega_{t-1}
	\label{eq:33:bop_solved} \\
	\bar{\kappa}_t(\omega_{t-1})
	&= \frac{1}{\nu} \left[ x_t - m_c c_k
	+ (m_c c_k - m_w) \omega_{t-1} \right]
	\label{eq:33:kappa_bar}
\end{align}
This equation formalizes the structuralist critique: the base
capacity to import capital goods is determined by export
earnings minus the baseline luxury consumption of the
capitalist class ($x_t - m_c c_k$). The sign of
$(m_c c_k - m_w)$ determines whether a rising wage share
expands or contracts the feasible mechanization gap.

\subsubsection{The Constrained Cost-Surplus Maximization
	Problem}

\paragraph{3.28}
The peripheral firm's tactical problem is therefore a
constrained cost-surplus maximization:
\begin{align}
	\max_{\kappa_t}\; c(\kappa_t)
	&= \Phi(\kappa_t) - \kappa_t \cdot MC(\omega_{t-1}, p^\kappa_t)
	\label{eq:33:constrained_obj} \\
	\text{subject to:}\quad & \kappa_t \le \bar{\kappa}_t
	\label{eq:33:constraint}
\end{align}
This formulation defines two distinct accumulation regimes,
endogenous to the interaction between domestic distributive
conflict and the global market.

\subsubsection{Endogenous Regime Shifts: The Corner Solution}

\paragraph{3.29}
The firm's unconstrained optimal technique
$\kappa_t^*(\omega_{t-1})$ is strictly increasing in the wage
share ($\partial \kappa_t^* / \partial \omega_{t-1} > 0$) due
to the Okishio labor-saving bias established in
\eqref{eq:32:cs_omega}.

\paragraph{3.30}
\textbf{Proposition 4.} \textit{Peripheral regime shifts.}
Two regimes are endogenous to the interaction between domestic
distributive conflict and the global market:
\begin{enumerate}
	\item \textbf{Regime A (Slack BoP):} When
	$\kappa_t^* \le \bar{\kappa}_t$, the firm solves the
	standard first-order condition, yielding $\kappa_t^*$. The
	labor-saving bias successfully extracts the Okishio
	distributive dividend.
	\item \textbf{Regime B (Peripheral realization trap):} When
	$\kappa_t^*(\omega_{t-1}) > \bar{\kappa}_t(\omega_{t-1})$,
	the constraint binds and the firm is forced into a corner
	solution:
	\begin{equation}
		\kappa_t = \bar{\kappa}_t
		\label{eq:33:corner}
	\end{equation}
\end{enumerate}
The corner solution arises through two intersecting structural
vulnerabilities. First, if $m_c c_k$ is large (the
conspicuous consumption drain), the baseline FX availability
is severely squeezed, lowering the intercept of
$\bar{\kappa}_t$. Second, if domestic agriculture is highly
inelastic ($m_w > m_c c_k$), a rising wage share
simultaneously increases the \textit{demand} for capital
imports (via $\kappa_t^*$) while \textit{contracting} the
supply of available FX (via $\bar{\kappa}_t$).
\textit{Proof.} Immediate from the Kuhn--Tucker conditions
applied to \eqref{eq:33:constrained_obj}--\eqref{eq:33:constraint}.
$\hfill\blacksquare$

\subsubsection{Effective Profit Compression and the Blocked
	Surplus}

\paragraph{3.31}
\textbf{Definition 3.} \textit{FX-Conversion Ratio.} The
FX-conversion ratio $\phi_t$ is the ratio of the maximum
feasible technique to the optimal unconstrained technique:
\begin{equation}
	\phi_t = \min\left(1, \frac{\bar{\kappa}_t}{\kappa_t^*}\right)
	\label{eq:33:phi}
\end{equation}
When the external constraint binds ($\kappa_t^* > \bar{\kappa}_t$),
$\phi_t < 1$. This generates a direct compression of the
effective profit share available for intensive accumulation.

\paragraph{3.32}
\textbf{Definition 4.} \textit{Blocked surplus.} Even if
domestic wage suppression yields a high nominal profit share
($\pi_t$), the physical inability to convert that surplus into
imported machinery squeezes the effective profit rate:
\begin{equation}
	\pi_t^{eff} = \pi_t \cdot \phi_t
	= (1 - \omega_{t-1}) \cdot \phi_t
	\label{eq:33:pi_eff}
\end{equation}
The difference constitutes the \textbf{blocked surplus}:
\begin{equation}
	\text{Blocked Surplus} = \pi_t(1 - \phi_t)
	\label{eq:33:blocked}
\end{equation}
Because the accumulation of machinery is driven by the
\textit{effective} profit share
($g_{K^{ME}} = \chi^{ME} \pi_t^{eff} - \delta^{ME}$), the
constraint decouples nominal profitability from capacity
formation. If capitalists squeeze wages to raise $\pi_t$, they
inadvertently increase conspicuous consumption ($m^C_t$),
which lowers $\bar{\kappa}_t$ and compresses $\phi_t$.

\subsubsection{The Regime-Shifting Aggregate
	Transformation Elasticity}

\paragraph{3.33}
\textbf{Definition 5.} \textit{BoP binding ratio.} The BoP
binding ratio $\tau_t$ is the ratio of productive capital
imports to import capacity:
\begin{equation}
	\tau_t \equiv \frac{\nu \cdot I^{ME}_t}
	{\text{ToT}_t \times X_t}
	\label{eq:33:tau}
\end{equation}
where $I^{ME}_t = K^{ME}_t - (1-\delta^{ME}) K^{ME}_{t-1}$
is gross investment in machinery and equipment, derived from
the perpetual inventory identity, and $\text{ToT}_t \times X_t$
proxies the import capacity (FX supply). Under Assumption~3,
productive capital imports $m^K_t = \nu \kappa_t$ are proxied
empirically by gross investment in machinery and equipment
$I^{ME}_t$, since mechanization is import-biased in the
periphery. The regime switch is governed by $\tau_t$:
\begin{itemize}
	\item $\tau_t < \bar{\tau}$: BoP constraint is
	\textbf{slack} (Regime A: unconstrained optimum).
	\item $\tau_t \ge \bar{\tau}$: BoP constraint is
	\textbf{binding} (Regime B: corner solution).
\end{itemize}
The threshold $\bar{\tau}$ is determined empirically by the
threshold regression in Section~4.

\paragraph{3.34}
The peripheral extension generates a regime-shifting
aggregate transformation elasticity. Under \textbf{Regime~A}
(slack BoP), the constraint does not bind, and the
elasticity coincides with the unconstrained expression from
Section~3.2:
\begin{equation}
	\theta^{A}_t = \theta^{SpecB}_t
	= \frac{\theta_{NRC}
		+ \left(\theta_\kappa + \theta_{\omega\kappa}\omega_{t-1}
		+ \theta_{\kappa p\kappa} p^\kappa_t\right)
		\frac{d\kappa_t^*}{dk^{NRC}_t}}
	{1 + s_{ME,t}\frac{d\kappa_t^*}{dk^{NRC}_t}}
	\label{eq:33:theta_A}
\end{equation}

\paragraph{3.35}
Under \textbf{Regime~B} (binding BoP), the firm is forced
into the corner solution $\kappa_t = \bar{\kappa}_t$. The
marginal capacity payoff at the corner is strictly positive
because the firm would like to mechanize further but cannot:
\begin{equation}
	MP_\kappa\big|_{\kappa=\bar{\kappa}}
	= \Phi'(\bar{\kappa}_t) - MC(\omega_{t-1}, p^\kappa_t) > 0
	\label{eq:33:MP_corner}
\end{equation}
The Kuhn--Tucker multiplier on \eqref{eq:33:constraint} is
strictly positive, confirming that the constraint binds. The
aggregate transformation elasticity under the binding regime
becomes:
\begin{equation}
	\theta^{B}_t
	= \frac{\theta_{NRC}
		+ MP_\kappa\big|_{\kappa=\bar{\kappa}}
		\cdot \frac{d\bar{\kappa}_t}{dk^{NRC}_t}}
	{1 + s_{ME,t}\frac{d\bar{\kappa}_t}{dk^{NRC}_t}}
	\label{eq:33:theta_B}
\end{equation}
Since $\bar{\kappa}_t$ is determined by the BoP constraint
\eqref{eq:33:kappa_bar} rather than by the first-order
condition \eqref{eq:32:foc_kappa}, the growth path
$d\bar{\kappa}_t / dk^{NRC}_t$ differs from the unconstrained
path $d\kappa_t^* / dk^{NRC}_t$. The FX-conversion ratio
$\phi_t = \bar{\kappa}_t / \kappa_t^* < 1$ compresses the
intensive margin, and the elasticity under the binding regime
is strictly lower:
\begin{equation}
	\theta^{B}_t < \theta^{A}_t
	\label{eq:33:theta_inequality}
\end{equation}

\paragraph{3.36}
\textbf{Proposition 5.} \textit{Regime-shifting elasticity.}
The peripheral aggregate transformation elasticity is a
regime-dependent parameter:
\begin{equation}
	\theta^{periphery}_t
	= \theta^{A}_t \cdot \mathbf{1}\{\kappa_t^* \le \bar{\kappa}_t\}
	+ \theta^{B}_t \cdot \mathbf{1}\{\kappa_t^* > \bar{\kappa}_t\}
	\label{eq:33:theta_regime}
\end{equation}
The elasticity shifts discretely when the BoP constraint
becomes binding. The regime switch is governed by the
interaction between the desired mechanization gap
$\kappa_t^*(\omega_{t-1}, p^\kappa_t)$ and the maximum
feasible gap $\bar{\kappa}_t(\omega_{t-1})$ determined by
export capacity, conspicuous consumption, and wage-goods
imports. An increase in the wage share $\omega_{t-1}$
simultaneously raises $\kappa_t^*$ (via the Okishio
labor-saving bias) and can contract $\bar{\kappa}_t$ (via
the wage-goods import channel when $m_w > m_c c_k$), making
the regime switch more likely. This is the peripheral
realization trap: distributive pressure that would induce
productive mechanization in the center instead triggers the
BoP constraint in the periphery, truncating the
transformation elasticity.

\subsubsection{The Failure of the Okishio Mechanism and
	Structural Inflation}

\paragraph{3.37}
\textbf{Proposition 6.} \textit{Blocked Okishio mechanism.}
At the boundary $\kappa_t = \bar{\kappa}_t$, the
implementation of the Okishio viability condition is blocked.
The firm satisfies the viability criterion---the technique
would reduce unit costs---but cannot deploy it because the
required machinery cannot be imported. The wage-savings
dividend cannot be capitalized into physical capacity. The
firm suffers a persistent shortfall in its cost-surplus
relative to its unconstrained optimal state.
\textit{Proof.} At $\kappa_t = \bar{\kappa}_t$, the
Kuhn--Tucker multiplier on \eqref{eq:33:constraint} is
strictly positive, so the first-order condition
$\partial c / \partial \kappa_t = 0$ does not hold. The
marginal viability yield exceeds the marginal cost at the
boundary, but the firm cannot expand further.
$\hfill\blacksquare$

\paragraph{3.38}
Unable to extract the distributive dividend via mechanization,
capitalists attempt to recover the blocked surplus through the
only remaining channel: price markups. When the wage-goods
bottleneck forces nominal wage adjustments and the external BoP
bottleneck blocks defensive mechanization, any further increase
in the wage share directly crushes profitability. The resulting
defensive pricing behavior ignites the structural inflation
dynamic \citep{Sunkel1958, Prebisch1981}. This is the
peripheral realization trap: nominal profitability coexists
with truncated mechanization and stalled structural envelopes,
permanently destabilizing the long-run coupling between plant
and machinery.


\subsection{Synthesis: What the Three Models Jointly
	Demonstrate}

\paragraph{3.38}
The three models developed in this section formalize a single
theoretical claim: that decentralized, individually
cost-minimizing decisions generate aggregate
disproportionality under the unbalanced growth closure
\citep{Okishio1961, Basu2026}.
Table~\ref{tab:coherence} summarizes the articulation
between the Okishio and Vidal frameworks that underpins
this claim.

\begin{table}[H]
	\centering
	\small
	\caption{Theoretical Coherence: Market Anarchy and the Internally Divided Firm Articulation Across Three Models}
	\label{tab:coherence}
	\begin{tabular}{p{2.2cm}p{3.8cm}p{3.8cm}p{4.2cm}}
		\toprule
		\textbf{Element} & \textbf{Descentralized Accumulation} &
		\textbf{Internally Divided Firm} & \textbf{Formal model captures} \\
		\midrule
		\midrule
		Micro logic & Viability condition (cost
		reduction at prevailing prices) & Internal
		contradiction (control vs.\ cooperation) &
		Optimization problem = cost-minimizing pole;
		constraints = cooperation limits \\
		Macro logic & Decentralized decisions $\to$
		disproportionality & Divided firm $\to$
		contested operating conditions & $\theta \neq 1$
		captures gap between accumulation and capacity \\
		Disequilibrium & Unbalanced growth (supply
		$\neq$ demand) & ``Mediocre sufficiency''
		\citep{Vidal2022} (no coherent optimum) & $\mu$
		demonstrates gap; $\theta$ captures structural
		inefficiency \\
		Peripheral specificity & BoP constraint blocks
		\textit{implementation} of viability & External
		constraint compounds internal division & Corner
		solution: $\kappa = \bar{\kappa}$ when
		constraint binds \\
		\bottomrule
	\end{tabular}
\end{table}

\paragraph{3.39}
The formal model does not assert equilibrium. It
demonstrates the negative: that capitalist accumulation,
left to its own decentralized logic, does not generate
proportional reproduction. The capacity utilization rate
$\mu$ is the empirical signature of this
disproportionality: it shows that the normal operating
conditions of capitalist firms, determined by their
cost-minimizing mechanization choices, do not meet
aggregate demand conditions. The transformation elasticity
$\theta$ captures this structural gap between accumulation
and realized capacity. Whether $\theta$ exceeds or falls
below unity is an empirical question determined by the
interaction of distributive and price channels under a
given accumulation regime. In the center, this gap is
cyclical---mediated by institutional configurations that
can be exhausted and reconstituted
\citep{Vidal2015, Vidal2019}. In the periphery, this gap
is permanent---locked in by the structural ceiling of
external subordination \citep{Prebisch1981, Ahumada2019}.

\paragraph{3.40}
Section~4 consolidates the empirical identification
strategy. It translates the three analytical models into
estimable cointegrating specifications, derives the
threshold regression for the peripheral corner solution,
and specifies the capacity recovery protocol. The
empirical analysis bounds this multi-sector mechanism to
a single-sector aggregate, which masks internal sectoral
divergences but preserves the core transformation logic.